\section{Results and Discussion}
\subsection{Current Evidence}
The repository contains a defined scope, acceptance criteria, GitHub issue and
milestone tracking, published Wiki pages, experiment templates, an agent log,
and presentation artifacts. These demonstrate project preparation, not the
correctness of the proposed automation. One owner-supplied OCR text sample has
been compared with the reference; runtime, successful form-generation counts,
and recovery results remain unmeasured.

On November 8, 2026, the owner reported that the sidebar setup worked and that
valid links were accepted while invalid links were rejected. This is a limited
owner-reported smoke check, not a complete usability or permission test.
A privacy-reviewed, avatar-free fixture and manual reference are available:
19 display rows and 17 expected students after staff exclusion, including two
in-person students. These counts describe test preparation, not successful
automated extraction. The next experiment compares raw OCR with the reference
before testing roster filtering and attendance ordering.

\subsection{Initial OCR Smoke Test}
On November 8, 2026, the owner reported that the integrated OCR output worked
as intended. This supports basic functionality in the tested environment but
does not establish general recognition accuracy. The owner subsequently provided
raw sidebar text containing 20 lines. After trimming surrounding whitespace and
aligning rows in sequence, 18 of 19 reference display rows matched exactly
(94.7\%). The remaining row, \texttt{Quinn Harmon - IP}, appeared as
\texttt{Quinn Harmon - I} followed by a standalone \texttt{IP}. This is one
incorrect row and one extra fragment, not a missing participant. The denominator
includes staff rows; this is not a character-accuracy metric.

The raw text is retained unchanged in the Cloud Vision sandbox. Exact deployed
revision, input hash, timing, and complete API response were not supplied.
Manual correction must reconstruct the suffix and remove the extra fragment;
blind concatenation would leave a duplicated letter. This motivates explicit
review before attendance classification or Form generation. No successful
filtering or Form generation is inferred.

\subsection{Review and Batch Progress}
The owner reported successful editable review/approval, the compact sidebar,
and three-image batch processing. The supplied approval summary contained 17
included entries, three excluded OCR rows, two in-person students, and 15 online
students. These are owner-reported checks, not independent usability results.
The latest change skips exact cross-screenshot overlaps; live confirmation
remains pending. Distinct same-name students may require manual addition.
Production Form generation is not yet validated; the separate feasibility trial
is described below.

\subsection{Google Integration Trial Status}
A separate two-student Google integration sandbox was implemented on November
13, 2026. It creates test copies and a roster, records verification observations,
and tracks outputs for cleanup. Seventy-two assertions passed using mocked
Google services in an injected JavaScript runtime after field-level template
diagnostics were added on November 15, 2026. The owner reported successful
copy creation, folder placement, naming, copied questions, and submissions after
manual publication. Counts progressed from one/zero to one/one, and each copy
contained only its intended answer. These are owner-reported observations;
respondent account categories and exact execution dates were not supplied.
The original remained unpublished with zero responses. After deploying the
field-level diagnostics, the owner reported all template checks passing, with
no review or failure results. No current tracked discrepancy was detected;
the cause of an earlier aggregate failure remains unconfirmed. Timestamp-only
differences require review rather than imply content changes; other differences
fail and read errors are explicit. These results support the core two-student
feasibility trial, not complete verification of quiz settings or access policy.
The owner subsequently confirmed submissions from signed-out users without a
sign-in prompt and from external personal accounts. The organization sets the
standard template policy; these observations match its current configuration
and supersede the earlier mandatory-sign-in assumption. Required name and email
answers support self-reported matching, not identity authentication. Resource
cleanup after the successful trial was subsequently confirmed by the owner:
only generated test copies and the roster were removed, with the template,
folder, control spreadsheet, and report retained. This does not establish live
failure or cancellation cleanup. Full settings fidelity remains unconfirmed. Full milestone acceptance
is separate from this result.
The experiment is not the production generation workflow.

\subsection{Control Spreadsheet Isolation}
The owner reported that edits and cancellation in review A did not affect
review B in a separate spreadsheet after following the approval-isolation test.
This is owner-reported cross-document isolation evidence, not validation of
first-use authorization by another account or concurrent use of one document.
The owner subsequently clarified that deployment uses one shared control
spreadsheet operated by student IT staff, with only one operator at a time.
The owner subsequently confirmed that canceling or finishing a review does not
carry state into the next review. This supports review reuse, not validation of
production-output preservation or simultaneous-user behavior.

\subsection{Continuous Integration and M2 Acceptance}
M2 closed November 15, 2026 with owner acceptance and an explicit exception for
hosted continuous integration (CI), blocked before execution by account billing.
On November 20, billing was resolved and successful runs were verified for the
\href{https://github.com/Michael-Kowata/ITC-492/actions/runs/43771430365}{checkpoint push},
\href{https://github.com/Michael-Kowata/ITC-492/actions/runs/43771457473}{pull request}, and
\href{https://github.com/Michael-Kowata/ITC-492/actions/runs/43773008874}{merged-main rerun}.
The exception is now resolved; the original timing remains historical evidence.
The workflow runs six JavaScript mock suites on Node 22 and checks whitespace.
Earlier local injected-runtime checks passed 239 assertions. Automated checks
do not exercise live Google services or establish browser usability; those are
separate from owner-reported trials. M2 acceptance does not imply completion
of production generation or recovery.

\subsection{Presentation Feedback}
The proposal presentation was scheduled for September 28, 2026. On October 28,
the project owner reported that it went successfully. The reported feedback was
to write the paper as the project develops and to organize it as a report across
multiple pages rather than leave it as a single-page outline. This revision
responds by expanding the proposal, separating major topics, and establishing
an ongoing documentation process. The actual delivery date and demonstrated
commit have not been independently confirmed in this draft.

\subsection{Limitations}
The sidebar now includes a server-side Cloud Vision request and a raw-text
result view. The single owner-supplied result is limited evidence; mocked
request/error checks cannot establish Google authorization or recognition quality.

Image quality, display-name ambiguity, service permissions, and execution limits
may constrain the solution. Human review reduces the effect of extraction errors
but introduces work that must be included in evaluation. A successful sample
does not establish support for every Zoom layout or account configuration.
